\documentclass[11pt,letterpaper]{article}
\usepackage[margin=0.88in]{geometry}
\usepackage[T1]{fontenc}
\usepackage{lmodern,microtype,amsmath,amssymb,amsthm,mathtools}
\usepackage{booktabs,longtable,array,enumitem,xcolor,hyperref,xurl,fancyhdr}
\hypersetup{colorlinks=true,linkcolor=blue!45!black,urlcolor=blue!45!black,citecolor=blue!45!black,
pdftitle={Saturating a Pell witness to remove a radical},
pdfauthor={OpenAI ChatGPT; prepared for Vladimir Reshetnikov},
pdfsubject={Explicit Diophantine equations equivalent to syntactic consistency of ZFC},
pdfkeywords={Diophantine equation, ZFC, consistency, Pell equation, coprime saturation}}
\setlength{\parindent}{0pt}\setlength{\parskip}{5pt}
\setlength{\emergencystretch}{2em}
\setlist{itemsep=3pt,topsep=5pt}
\newtheorem{theorem}{Theorem}[section]
\newtheorem{lemma}[theorem]{Lemma}
\newtheorem{proposition}[theorem]{Proposition}
\newtheorem{corollary}[theorem]{Corollary}
\theoremstyle{remark}\newtheorem{remark}[theorem]{Remark}
\newcommand{\Z}{\mathbb Z}\newcommand{\N}{\mathbb N}
\newcommand{\Con}{\operatorname{Con}}\newcommand{\ZFC}{\mathrm{ZFC}}
\newcommand{\fl}[1]{\left\lfloor#1\right\rfloor}\newcommand{\rg}[2]{\{#1,\ldots,#2\}}
\newcommand{\sha}{\texttt{84e16e43092614bae262ca59651cb9d77e31a255}}
\pagestyle{fancy}\fancyhf{}\fancyhead[L]{\small Saturated Pell consistency equations}
\fancyhead[R]{\small 24 October 2026}\fancyfoot[C]{\thepage}
\setlength{\headheight}{14pt}
\title{Saturating a Pell witness to remove a radical\\[4pt]
\large Smaller explicit Diophantine equations for syntactic consistency of ZFC}
\author{Research contribution prepared for Vladimir Reshetnikov\\
\small AI-assisted mathematical analysis and artifact preparation: OpenAI ChatGPT}
\date{24 October 2026}
\begin{document}
\maketitle
\begin{abstract}
We give four explicit polynomials in ten strictly positive integer
unknowns, with operation/degree profiles
$92/283$, $93/278$, $98/258$, and $99/254$.
They improve the inspected ProveIt ten-positive baselines $99/368$ and
$105/340$.  The substantive change absorbs the first Pell ordinate into
the strong Pell witness, making two required squares coprime factors of
one square.  An elementary congruence excludes the exceptional cancellation
between the two remaining radicals.  The final equation is consequently
a four-factor norm instead of an eight-factor norm.  Both directions of
positive-witness existence are proved; the reverse direction uses a finite
Pell-matrix period and the explicit auxiliary witness already available in
the source construction.

The same saturation principle, with a compatible rescaling of the
auxiliary coordinate and full-residue interpolation, gives an equation
in twelve unrestricted integers with profile $123/1084$.  It improves
both current twelve-integer parents.  Polynomial coordinate changes
with profiles $143/1200$ and $146/1184$ are retained as controls in an
appendix.  A final appendix develops an
explicit degree-14 local equational proof kernel as a component of a
different future source encoding.

The consistency calibration is inherited from the precisely identified
ordinary native, compiler, and refutation theorems in the repository.  This
contribution supplies ordinary proofs formalizable in ZFC and reproducible
finite certificates; it does not supply a new end-to-end Lean or Rocq proof,
prove consistency, establish a global optimum, or claim a new minimum
number of unrestricted integer variables.
\end{abstract}
\tableofcontents
\newpage

\section{Results, comparison, and the logical claim}

The source snapshot is ProveIt commit \sha.  All comparisons in this
article refer to that snapshot.  Its active research directory is more
recent than its placed-manuscript catalogue: the complete displayed page
has $101$ operations and degree $368$, while the accepted zero-offset
evaluator already has $99$ operations at the same degree.  The direct
coordinate sibling has $105$ operations and degree $340$.
The new comparison is therefore against these later accepted evaluators,
not against the older wave-ten catalogue \cite{repo,zero,direct}.

\begin{center}
\begin{tabular}{@{}llrr@{}}\toprule
Construction & Quantified domain & Binary operations & Exact degree\\\midrule
Pinned zero-offset parent & $10$ positive integers & $99$ & $368$\\
New saturated upper center & $10$ positive integers & $92$ & $283$\\
New saturated sharp center & $10$ positive integers & $93$ & $278$\\
Pinned direct parent & $10$ positive integers & $105$ & $340$\\
New direct upper center & $10$ positive integers & $98$ & $258$\\
New direct sharp center & $10$ positive integers & $99$ & $254$\\\midrule
Pinned rational scalar carrier & $12$ unrestricted integers & $145$ & $1200$\\
Pinned integral scalar carrier & $12$ unrestricted integers & $141$ & $1456$\\
New saturated scalar carrier & $12$ unrestricted integers & $123$ & $1084$\\
Centered integral control & $12$ unrestricted integers & $143$ & $1200$\\
Centered rational control & $12$ unrestricted integers & $146$ & $1184$\\\bottomrule
\end{tabular}
\end{center}

An operation is one binary addition, subtraction, or multiplication,
including multiplication by a numerical constant and squaring.  Registers
can be shared.  Constants and variable inputs are leaves, not operations.
Degree counts quantified variables only.  This is an arithmetic-circuit
measure, not bit complexity: the large fixed coefficients are reported
separately and are not being made smaller.

The $92/283$ form saves seven operations and $85$ degrees from $99/368$,
approximately $7.1\%$ and $23.1\%$, respectively.  The $99/254$ direct form
uses no more operations than the former cheapest ten-positive evaluator
and reduces its degree by $114$, approximately $31.0\%$.
The four new ten-positive points are different cost/degree choices.
The twelve-integer saturation saves $22$ operations and $116$ degrees
from its rational parent, and $18$ operations and $372$ degrees from
its integral parent.
They do not improve the repository's nine-positive or ten-unrestricted-
integer variable counts, nor its separate low-degree quartic families.

\begin{theorem}[Consistency interpretation and its dependency boundary]
Let $N_{2,\varepsilon}$ and $N^{\rm dir}_{2,\varepsilon}$ be the completely
specified polynomials below, for $\varepsilon=0,1$, at the coherent fixed
tuple of Appendix~\ref{sec:constants}.  The ordinary source calibration
identified in Section~\ref{sec:dependencies}, composed with the witness
constructions in this article, gives
\[
 \ZFC\vdash
 \left[\neg\exists\mathbf x\in\Z_{>0}^{10}\,
             N_{2,\varepsilon}(\mathbf x)=0\right]
       \ \longleftrightarrow\ \Con(\ZFC),
\]
and the same statement for the direct variant.  The twelve-integer
form in Section~\ref{sec:integer12-saturated}, and the controls in
Appendix~\ref{sec:integer12-main}, have the analogous statement with
$\mathbf x\in\Z^{12}$.
\end{theorem}

Here $\Con(\ZFC)$ denotes absence of a finite syntactic refutation in the
same effective ZFC calculus used by the pinned source.  The displayed
claim is an ordinary mathematical derivability claim, using the existing
ordinary calibration, formalizable in ZFC.  The archive does not contain an
independent formal verification of that entire inherited theorem.  Its
new machine-executed checks concern the literal polynomial specifications,
their algebraic identities, and their finite certificates.  No assertion
that either side of the equivalence is true is needed for the construction.

\subsection{What is new in this round}

The transport expansion $W B/(q-1)=LB+B/(q-1)$ and the general ideas of
centering and saturating Pell witnesses already occur in the source
collection.  They are credited as such.  The principal new result is the
complete application of coprime-product saturation to the current
ten-positive kernel: a two-root recovery proof, an inexpensive upper
center, full positive completeness after changing the strong witness,
and the four literal equations above.  A focused search of the downloaded
incoming submissions and the inspected active proofs found no earlier
instance of these complete equations.  This is a repository-relative
comparison, not a claim of worldwide priority for coprime-square splitting.
The unrestricted-integer extension applies the same idea to the signed
scalar carrier.  Its additional contribution is the compatible auxiliary
rescaling and residue interpolation needed to preserve completeness.

\section{The complete ten-positive polynomial}
\label{sec:definition}

Choose once and for all $\varepsilon\in\{0,1\}$.  It selects a polynomial;
it is not an eleventh unknown.  The ten unknowns are
\[
 J,F,\alpha,i,T,w,Y,n,\eta,\zeta\in\Z_{>0}.
\]
The four fixed integers $\beta,K_0,M_C,M_F$ are specified completely in
Appendix~\ref{sec:constants} and as literal payloads in the archive.
They satisfy
\[
 \beta\ge31\text{ odd},\qquad K_0>0,\qquad
       0<M_C<\beta<M_F<2\beta.
\]
All symbols introduced next are polynomial abbreviations; none adds an
existential variable.  Set
\begin{align}
 q&=\beta J+1,& C&=q-F-\alpha,& X&=wq,\nonumber\\
 \mu&=q^2+1,& k&=\eta+\zeta,& c&=kY+\eta,\label{eq:outer}\\
 A&=\mu Y(X+1)+1,& H&=4A-5,&\Delta&=A^2-1,\nonumber\\
 V_0&=X\mu^2Y^2,& E&=1+V_0(V_0+1)k^2,&U&=ic^2.\label{eq:kernel}
\end{align}
The rank and transport data are
\begin{align}
 R&=[(q-1)(q-F)+\alpha](q^2-1)+(M_C+qM_F)J,\nonumber\\
 e&=k-R-1,& B_t&=(K_0+w)C-F,\nonumber\\
 T'&=T+\alpha,& b&=cT'-R,&L&=kT'.\label{eq:transport}
\end{align}
The saturated strong and auxiliary quantities are
\begin{align}
 D_f&=1+\Delta U^2E,&Q&=\Delta^2U^2E,&\delta&=Q-1,\nonumber\\
 V&=D_fb-c,&D_y&=QV^2-1,&D_g&=ED_f,\nonumber\\
 C_g&=(A+\varepsilon)UE,&K_+&=X(3c-2X)+2.\label{eq:saturated}
\end{align}
Define the integer polynomials
\begin{align}
 P_1&=q^3(q-1),&M&=P_1\mu^2H,\nonumber\\
 Z&=P_1\mu^2\{H(n+V+C_g-LB_t)+K_+\}\nonumber\\
   &\hspace{13mm}+H\{\mu^2[(q-1)Y-q^3B_t]-P_1e\},\label{eq:numerator}\\
 f_s&=\delta Z^2,& f_g&=\delta M^2D_g,&f_y&=M^2D_y.\nonumber
\end{align}
The promised single equation is
\begin{equation}
 \boxed{\quad N_{2,\varepsilon}
       =(f_s-f_g-f_y)^2-4f_gf_y=0.\quad}
 \label{eq:main}
\end{equation}
There is no division, radical, variable exponent, unbounded product, or
undetermined program parameter in this polynomial.  Large powers used in
the fixed coefficient recipe define constants before the unknowns are
introduced.  The names \texttt{saturated\_upper.json} and
\texttt{saturated\_sharp.json} correspond to $\varepsilon=1$ and $0$.

\subsection{The direct coordinate version}

Replace the unknown $w$ by a positive unknown $x$ and set $X=q+x$.
In \eqref{eq:transport} replace $B_t$ by
\[
 B=(qK_0+X)C-qF.
\]
Use $P_1=q^4(q-1)$ and $M=P_1\mu^2H$.  Replace only the numerator by
\begin{align}
 Z_{\rm dir}
 &=P_1\mu^2\{H(n+V+C_g-LB)+K_+\}\nonumber\\
 &\quad+H\{\mu^2[q(q-1)Y-q^3B-(q-1)X]-P_1e\}.
 \label{eq:directnumerator}
\end{align}
All other formulas, including the final norm, stay as displayed.
In particular the direct $P_1$ contains $q^4$, not $q^3$.
Its normalized center has the additional term $-X/q^4$; this is a proof
identity, not a division in the evaluator.  The retained direct-coordinate
argument recovers $q\mid X$ only after the whole-product projection has
removed the $H$ denominator.  No earlier cancellation of $q$ is used.

\subsection{The exact arithmetic schedules}

The JSON files list every gate in topological order as
\texttt{[register, operation, operand1, operand2]}.  All gates, inputs,
and fixed ports are live.  For efficient strong arithmetic, put
\[
 U_E=UE,\quad \Delta=(A-1)(A+1),\quad
 S_E=\Delta U_E,\quad B_f=S_EU,
 \quad D_f=B_f+1,\quad Q=\Delta B_f.
\]
For the upper center $(A+1)U_E$, the literal $A+1$ is already needed in
$\Delta$.  The sharper $AU_E$ requires one extra addition to obtain $A$.
The final norm takes thirteen gates from $M^2$ through its output,
including the product $D_g=ED_f$.  The operation breakdowns are
\begin{center}
\begin{tabular}{@{}lrrrr@{}}\toprule
Variant & Multiplications & Additions & Subtractions & Total\\\midrule
Upper &55&20&17&92\\
Sharp &55&21&17&93\\
Direct upper &59&21&18&98\\
Direct sharp &59&22&18&99\\\bottomrule
\end{tabular}
\end{center}
The distributed transport numerator used here was already available in
the earlier dyadic submission.  The new schedules compose it with the
saturated kernel, not with a re-labelled old finalizer.

\section{Why the product of two squares can replace two radicals}
\label{sec:product-explanation}

The central design choice is the congruence
\[
 D_f=1+\Delta U^2E\equiv1\pmod E.
\]
Thus $E$ and $D_f$ are coprime.  Once the two-root norm forces
$ED_f$ to be a square, both factors must individually be squares:
\[
 E=\tau^2,\qquad D_f=f^2.
\]
The proof-only integer $I=U\tau$ then satisfies the original strong
equation
\[
 f^2-\Delta I^2=1,\qquad c^2\mid I.
\]
There is no additional unknown for $I$ or $\tau$ in the new equation.

One must also exclude a false zero caused by cancellation of two equal
irrational roots.  Here that exclusion is built into the same factor $E$:
if
\[
 ED_f=\frac{QV^2-1}{Q-1},
\]
then $(Q-1)ED_f=QV^2-1$, whereas $E\mid Q$ makes the two sides
congruent to $0$ and $-1$ modulo $E>1$.  This contradiction is the reason
the two-root finalizer is sound.  No unsupported assertion that a rational
sum of radicals makes every radical rational is used.

The converse is also substantive.  An old strong ordinate need not be
divisible by $\tau$.  Instead, after constructing the canonical outer
witnesses, choose a positive Pell-matrix period modulo $c^2\tau$.
The resulting new ordinate is divisible by $c^2\tau$, so
$U=I/\tau=ic^2$ for a positive integer $i$.  The explicit auxiliary
Pell solution is large enough to make the final $n$ positive despite the
new upper center.  These details are proved in the next section.

\section{An unbounded arithmetic proof of the saturation step}
\label{sec:saturation-semantics}

We prove the equivalence for both charts and for either of the two
fixed centers $C_g=(A+\lambda)UE$, where $\lambda=\varepsilon\in\{0,1\}$.
Write $K_{\rm pos}=K_+$ and $G=D_g=ED_f$.  The other notation is that of the complete
equation above.  Throughout this section,
all ten independent coordinates are strictly positive integers.  A
``root'' denotes a positive real square root only inside the proof;
the polynomial evaluator has no root operation.

The inherited input is the accepted ten-positive zero-offset
construction, together with its coherent native/controller/refutation
calibration.  Its proof imports the universal Pell facts and the
source61 lemmas named \texttt{lem:rank}, \texttt{lem:index}, and
\texttt{lem:auxcomplete}.  We state their interfaces at the points
where they are used.  The old strong witness is not retained literally:
the proof will recover $I=U\sqrt E$ and, in the converse direction,
choose a new strong Pell period.  The complete original port tuple
and the complete BFS tuple are separate allowed calibrations;
individual ports are never interchanged.

\subsection{Two roots recover the three required Pell conditions}

\begin{lemma}[Root recovery by a coprime product]
\label{lem:saturation-root-recovery}
At every positive integer zero of the displayed polynomial, there
are positive integers $\tau,f,y$ such that
\begin{equation}
 E=\tau^2,\qquad D_f=f^2,\qquad
 QV^2-(Q-1)y^2=1,
 \label{eq:sat-three-identities}
\end{equation}
and
\begin{equation}
 \frac ZM=\varepsilon_g f\tau+\varepsilon_y y,
 \qquad \varepsilon_g,\varepsilon_y\in\{-1,1\}.
 \label{eq:sat-signed-factor}
\end{equation}
No sign restriction on $Z$ or $V$ is needed for this conclusion.
\end{lemma}

\begin{proof}
Before using the equation, $E>1$, $D_f>c+1$, $Q>1$, $M>0$, and
$\delta=Q-1>0$.  Since $V=D_fb-c$ and $b$ is an integer, either
$b=0$ and $V=-c$, or $|V|\ge D_f-c>1$.  In particular $D_y>0$.
Set
\[
 G=ED_f,\qquad Y_* =\frac{D_y}{\delta}
       =\frac{QV^2-1}{Q-1}.
\]
The two positive rational numbers $G,Y_*$ cannot coincide.  Indeed,
$G=Y_*$ would imply
\[
 Q(V^2-G)=1-G.
\]
Both $Q$ and $G$ are divisible by $E$, so this would give $E\mid1$.
This simple obstruction replaces a radical size or sign separation.

At a zero put $z=Z/M$.  The equation is exactly
\[
 (z^2-G-Y_*)^2=4GY_*.
\]
If $z=0$, then $G=Y_*$, which was excluded.  For $z\ne0$,
\[
 \left(\frac{z^2+G-Y_*}{2z}\right)^2=G.
\]
The complementary rational number has square $Y_*$.  Thus both
positive roots are rational.  Since $G$ is an integer, it is an
integer square.  Moreover
\[
 D_f=1+\Delta U^2E\quad\Longrightarrow\quad\gcd(E,D_f)=1.
\]
A square product of positive coprime integers has square factors.
Write $E=\tau^2$, $D_f=f^2$, and $g=f\tau=\sqrt G$.

For clarity, rational integrality of the remaining adjacent-Pell
root has the following elementary descent.  If $y=a/m>0$ in lowest
terms satisfies $Qv^2-(Q-1)y^2=1$ with integer $v>0$, then
$m^2\mid Q-1$.  For $v>1$ use
\[
 v'=(2Q-1)v-2(Q-1)y,
 \qquad y'=(2Q-1)y-2Qv.
\]
The first coordinate is integral, the same norm is preserved, and
$v<y<Qv/(Q-1)$ gives $-v<v'<v$.  The new first coordinate cannot
vanish.  After taking absolute values, its positive integer size
has decreased.  The reduced denominator of $y'$ is still $m$:
$2Q-1\equiv1\pmod m$ and $a$ is coprime to $m$.  Descent reaches
$v=1$, where $y=1$, so $m=1$.  Apply this argument initially to
$v=|V|$.  The signed factorization now gives
\eqref{eq:sat-signed-factor}.
\end{proof}

\subsection{The old strong witness and the early inequalities}

Define only inside the proof
\begin{equation}
 I=U\tau.
 \label{eq:sat-proof-I}
\end{equation}
Then $I$ is a positive integer, $c^2\mid I$, and
\begin{equation}
 f^2-\Delta I^2=1,\qquad Q=(\Delta I)^2,
 \qquad V=f^2(cT'-R)-c.
 \label{eq:sat-restored-kernels}
\end{equation}
These are exactly the normalized strong and auxiliary identities
of the accepted zero-offset proof.  Its old strong input could be
$i_{\rm old}=I/c^2=i\tau$; it is not the new literal coordinate $i$.

\begin{lemma}[Folding and size before transport]
\label{lem:sat-early-fold}
At every positive zero there are indices $p,n_0\ge1$, with $p$ odd,
such that
\[
 c=\psi_A(p),\qquad
 k=2\psi_{2V_0+1}(n_0),\qquad n_0\ge2.
\]
Consequently
\begin{equation}
 k\ge8V_0+4,
 \qquad c>kY\ge k>8q^5,
 \qquad K_{\rm pos}>0,
 \qquad H>3X>0.
 \label{eq:sat-early-sizes}
\end{equation}
These conclusions do not use the sign of the transport numerator.
\end{lemma}

\begin{proof}
The inherited strong/auxiliary folding lemma is uniform for $A\ge3$.
Its hypotheses are the strong norm, $c\le I$, $c^2\mid I$, the
adjacent auxiliary norm, and $V\equiv-c\pmod f$.  They follow from
\eqref{eq:sat-three-identities}--\eqref{eq:sat-restored-kernels}.
It gives $c=\psi_A(p)$ for odd $p\ge1$.  The unchanged first norm
gives the stated expression for $k$.  If $n_0=1$, then $k=2$;
positivity forces $\eta=\zeta=1$, and $c=2Y+1$ lies strictly
between $\psi_A(1)=1$ and $\psi_A(2)=2A$.  This is impossible.
Thus $k\ge2\psi_{2V_0+1}(2)=8V_0+4$.
The inequalities in \eqref{eq:sat-early-sizes} now use
$V_0=X\mu^2Y^2$, $X\ge q$, and $\mu=q^2+1$.  In particular
$c>2X$, so $K_{\rm pos}=X(3c-2X)+2>0$; the formula
$H=4\mu Y(X+1)-1$ gives $H>3X$.
\end{proof}

The port bounds of the native mask give, for every positive assignment,
\begin{align}
 R&<q^4+(\alpha+2)q^2,\notag\\
 |R|&<q^4+Fq^3+(\alpha+2)q^2,\notag\\
 \frac{|R|}{\mu^2}&<F+\alpha+1.
 \label{eq:sat-mask-bounds}
\end{align}
For the last inequality, multiplying by $\mu^2=(q^2+1)^2$ leaves
the positive difference
\[
 F(q^4-q^3+2q^2+1)+\alpha(q^4+q^2+1)+1.
\]
The first bound and $c>8q^5$ show that
$b=c(T+\alpha)-R>0$, hence $b\ge1$ and $V=D_fb-c>1$.
Put
\[
 u=y-V>0,\qquad \rho=\frac{e}{\mu^2}.
\]
The adjacent norm and $\Delta D_f=Q+\Delta<2Q-2$ give
\[
 u=\frac{V^2-1}{(Q-1)(y+V)}
   <\frac{V}{2(Q-1)}<\frac b\Delta.
\]
Using $c<2kY$, $\Delta>8Y$, $\Delta>\mu^2$, and
\eqref{eq:sat-mask-bounds}, we obtain
\begin{equation}
 u<\frac{kT'}4+F+\alpha+1,
 \qquad \rho<\frac{k}{\mu^2}+F+\alpha+1.
 \label{eq:sat-general-u-rho}
\end{equation}
If $C=q-F-\alpha\ge1$, then
$q(q-F)-C=q\alpha+(q-1)C>0$ and therefore $R>0$.
In that case the sharper inequalities are
\begin{equation}
 u<\frac{kT'}4,
 \qquad \rho<\frac{k}{\mu^2}.
 \label{eq:sat-sharp-u-rho}
\end{equation}
At this stage $\rho$ may still be rational.  No denominator has yet
been cancelled.

\subsection{Why the product center suffices on both charts}

\begin{lemma}[Upper and sharp centers]
\label{lem:sat-centers}
At every positive zero,
\begin{equation}
 g<C_g,
 \qquad 0<C_g-g<(\lambda+1)UE\le2UE.
 \label{eq:sat-center-loss}
\end{equation}
For the sharp center $C_0=AUE$, the stronger estimate
\begin{equation}
 C_0-g>\frac{UE}{4A}>X
 \label{eq:sat-direct-gap}
\end{equation}
also holds.  In particular, for either center and either root sign,
$\varepsilon_g g-C_g+X/q^4<0$.
\end{lemma}

\begin{proof}
Since $I>1$ and $f^2=A^2I^2-I^2+1$,
$(A-1)I<f<AI$.  Multiplication by $\tau$ proves
\eqref{eq:sat-center-loss}.  For the sharper lower bound use
\[
 C_0^2-G=U^2E^2-E.
\]
Here $U^2E>2$ and $g<C_0$, so
$(C_0-g)>UE/(4A)$.  The next estimate is unconditional on the
positive coordinate domain: $k\ge2$, $c>kY$, $U\ge c^2$, and
$E>V_0^2k^2$ imply
\[
 UE>V_0^2k^4Y^2\ge16V_0^2Y^2.
\]
Also $A=\mu Y(X+1)+1<4\mu XY$.  Hence
\[
 \frac{UE}{4A}>
 \frac{V_0^2Y}{\mu X}=X\mu^3Y^5>X.
\]
In particular the direct estimate uses neither $q\mid X$ nor any
decoded source information.
\end{proof}

\subsection{Transport soundness without a large majorant}

On the ordinary chart write $W=1+(q-1)L$ and
\[
 t=\frac{B_t}{q-1},\qquad
 s=\frac{Y}{q^3}>0,\qquad
 \kappa=\frac{K_{\rm pos}}H>0.
\]
The exact normalized carrier and
\eqref{eq:sat-signed-factor} give
\[
 n=(\varepsilon_g g-C_g)+(\varepsilon_y y-V)
       -\kappa-s+Wt+\rho.
\]
By Lemma~\ref{lem:sat-centers}, every zero satisfies
\begin{equation}
 n<u+Wt+\rho.
 \label{eq:sat-transport-upper}
\end{equation}
Suppose $B_t\le-1$.  If $C\ge1$, then $Wt<-L=-kT'$;
\eqref{eq:sat-sharp-u-rho} gives
\[
 n<\frac{kT'}4+\frac{k}{\mu^2}-kT'<0.
\]
If $C\le0$, then $B_t=(K_0+w)C-F\le-F$ and $Wt<-FL$.
By \eqref{eq:sat-general-u-rho},
\[
 n<\frac{kT'}4+\frac{k}{\mu^2}
      +2F+2\alpha+2-FkT'<0.
\]
For completeness, the negative of the expression on the right is
increasing in $F,T,\alpha\ge1$ when $k\ge8$; its minimum is
$3k/2-k/\mu^2-6>0$.  Both cases contradict $n>0$.
We have proved $B_t\ge0$, and then
$(K_0+w)C\ge F>0$ gives $C>0$.

On the direct chart put $W=1+q(q-1)L$.  The normalized carrier differs by the transport
term $-WB/[q(q-1)]$ and the additional term $-X/q^4$.
Lemma~\ref{lem:sat-centers} absorbs the resulting $+X/q^4$ in the
equation for $n$.  Thus the same upper bound holds with
$Wt$ replaced by $WB/[q(q-1)]$.  If $B\le-1$ and $C\ge1$,
this term is smaller than $-L$.  If $C\le0$, then
$B=(qK_0+X)C-qF\le-qF$, so it is smaller than
$-qFL\le-FL$.  The identical two contradictions give $B\ge0$
and $C>0$.  Neither argument assumes the divisibility $q\mid X$.

\subsection{Licenses for the inherited index and projection proof}

We record the order of implications because it prevents circular
cancellation arguments.  Once transport is nonnegative, the native
mask estimate gives $3q<R<q^4-2$.  The recovered strong witness
satisfies $c^2\mid I$ and $c>2R$.  Since $c=\psi_A(p)>2$ and
$A\ge3$, one also has $c>2p$.

The source61 rank lemma says that $c=\psi_A(p)>2$,
$f^2-\Delta I^2=1$, $f,I>0$, and $c^2\mid I$ imply
\[
 f=\chi_A(m),\quad I=\psi_A(m),\quad
 pc\mid m,\quad c\mid m,\quad m>2p.
\]
Put $S=\Delta I$.  Then $S\ge2$, $c\mid S$, and
$S^2=\Delta(f^2-1)$.  The two required auxiliary congruences are
\[
 V\equiv-c\pmod f,
 \qquad V\equiv-R\pmod c;
\]
the second uses $f^2\equiv1\pmod c$.  The source61 index lemma,
with $0<p,R<c/2$, now gives $p=R$.

The index $R$ is odd.  If $q$ were odd, then $J$ would be even and
the native mask formula would make $R$ even.  Thus $q$ is even
and $J$ is odd.  Only now may one use
\[
 \gcd\bigl(\mu^2,q^3(q-1)H\bigr)=1
 \quad\text{or, on the direct chart,}\quad
 \gcd\bigl(\mu^2,q^4(q-1)H\bigr)=1.
\]
Indeed $H\equiv-1\pmod\mu$, and even $q$ makes
$\gcd(q-1,\mu)=1$.  The signed carrier is integral, so the
unchanged denominator argument gives $\mu^2\mid e$.

All following uses of the zero-offset parent are now licensed:
its first-index argument gives $n_0=(R+1)/2$; its offset-free
ratio estimate gives the same half-binomial interval and size bound;
and its whole-product projection gives $X=2^R$.  More explicitly,
on the ordinary chart the integral carrier yields
$H\mid q^3(q-1)K_{\rm pos}$.  Writing
\[
 K_{\rm quad}=X\bigl(X+2(A-2)c\bigr)-1,
 \qquad K_{\rm pos}=HXc-2K_{\rm quad},
\]
one cancels only $2$ (since $H$ is odd).  The universal projection
therefore yields
\[
 H\mid q^3(q-1)(X-2^R)(2^RX+1).
\]
The entire product has absolute value less than $H$ by the retained
ratio and height bounds; hence $X=2^R$.  One does not cancel an
unverified unit modulo $H$.  The direct parent supplies exactly
the corresponding product argument with its extra carrier factor.

The new center causes no change to these interfaces.  Both $C_g$
and $g=f\tau$ are integers, so the only nonintegral terms remain
the old fractions.  On the direct chart, after the main projection
makes $K_{\rm pos}/H$ integral, the residual relation is
\[
 \frac{Y}{q^3}-\frac{WB}{q(q-1)}-\frac{X}{q^4}\in\mathbb Z.
\]
Its numerator modulo $q$ gives $q\mid X$.  Since $X=2^R$,
$q$ is a power of two; $R>3q$ then gives $q^4\mid X$.
Set $w=X/q$ and use $B=qB_t$ to finish the old source decoding.

The remaining parent steps give $R\equiv3\pmod4$, the exact
half-binomial value of $Y$, and all native source divisibilities.
Its final reconstruction uses a \emph{fresh unscaled} source61
parameter $A_{\rm source}=Y(X+1)+2$.  It neither divides the new
strong witness nor identifies this unscaled parameter with $A$.
The fixed native/controller calibration then produces a finite
syntactic ZFC refutation.  This proves soundness for each new
center and chart.

\subsection{Completeness: choose a divisible Pell period}

Suppose a finite ZFC refutation exists.  The parent's canonical
construction gives the outer and first/main data, with
\[
 E=\tau^2,\quad c=\psi_A(R),\quad
 R\equiv3\pmod4,\quad R>3q,
\]
and canonical integers $t\ge0$, $\rho\ge0$, $s>0$, and
$\kappa>0$.  Choose the strong pair anew rather than reusing the
parent's particular index.

Let $N_0=c^2\tau$.  The determinant-one matrix
\[
 M_A=\begin{pmatrix}A&\Delta\\1&A\end{pmatrix}
\]
has finite positive order in the group of invertible matrices over
$\mathbb Z/N_0\mathbb Z$.  At a positive multiple $m$ of that order,
\[
 I=\psi_A(m)>0,\qquad f=\chi_A(m)>0,
 \qquad N_0\mid I,\qquad f^2-\Delta I^2=1.
\]
Set $i=I/(c^2\tau)>0$.  Then the literal $U=ic^2$ satisfies
$U=I/\tau$, and the new $D_f,Q$ are exactly $f^2,(\Delta I)^2$.
This construction requires divisibility of a Pell ordinate by a
prescribed integer, which a finite period supplies.  It requires
no intersection theorem for two unrelated Pell sequences.

The source61 auxiliary-completeness lemma assumes only $A\ge2$,
$R\ge3$, $R\equiv3\pmod4$, $c=\psi_A(R)$,
$f^2-\Delta I^2=1$, $f,I>0$, and $c\mid I$.  All hold for this
new pair.  In particular the lemma fixes neither the strong index
nor the parity of $I/c^2$.  Its explicit witness, with
$S=\Delta I$, $Q=S^2$, and $R=2v+1$, is
\begin{align*}
 V_*&=\chi_S(R)/S=Q_v(S^2),\qquad y=\psi_S(R),\\
 T_*&=\frac{V_*+c+Rf^2}{cf^2},\qquad
 QV_*^2-(Q-1)y^2=1.
\end{align*}
The lemma makes $T_*$ a positive integer.  Its usual Pell growth
bound gives, since canonical $v\ge2$,
\[
 V_*\ge(4Q-3)^v\ge(3Q)^v\ge9Q^2.
\]
Because $c<S$, $f^2<Q$, and $\alpha<q<S$,
\[
 T_*>\frac{V_*}{cf^2}>\frac{V_*}{S^3}\ge9S>\alpha.
\]
Thus $T=T_*-\alpha$ is positive and the literal auxiliary register
equals $V_*$.  For $u=y-V_*>0$, using $V_*\ge2$ and $y<2V_*$,
\begin{equation}
 u=\frac{V_*^2-1}{(Q-1)(y+V_*)}
   >\frac{V_*}{4Q}\ge\frac94S^2
   =\frac94\Delta^2I^2.
 \label{eq:sat-complete-gap}
\end{equation}

Here $U=ic^2>1$, so $UE=I\tau<I^2$.  Also
$\kappa<c<I$ and $s<Y<c<I$; the first bound follows from
$K_{\rm pos}<3Xc$ and $H>3X$.  Define the integer
\[
 n=(g-C_g)+(y-V_*)-\kappa-s+Wt+\rho.
\]
Its total negative loss is less than
\[
 2UE+\kappa+s<2I^2+2I<4I^2,
\]
whereas \eqref{eq:sat-complete-gap} exceeds $4I^2$, since
$\Delta\ge3$.  Thus $n>0$.  The normalized center equals
$g+y$, and one factor of the two-root norm vanishes.  This is
completeness on the ordinary chart for both centers.  The explicit
auxiliary witness provides the needed growth; no unproved
cofinality assertion has been inserted.

For the direct chart, transfer this newly constructed ordinary zero.
It has $X=2^R>q$, $q^4\mid X$, and $t=B_t/(q-1)\ge0$.
Set
\begin{align*}
 x&=X-q,\\
 W_{\rm ord}&=1+(q-1)L,\qquad
 W_{\rm dir}=1+q(q-1)L,\\
 n_{\rm dir}&=n_{\rm ord}
       +(W_{\rm dir}-W_{\rm ord})t+X/q^4.
\end{align*}
Every coordinate is positive.  All kernels and $C_g$ agree, and
the normalized centers agree.  In fact
$M_{\rm dir}=qM_{\rm ord}$ and $Z_{\rm dir}=qZ_{\rm ord}$,
so the new norm scales by $q^4$ and remains zero.

\begin{theorem}[Semantic equivalence of the four saturated equations]
\label{thm:saturation-equivalence}
For either fixed center and either chart, with any one complete
coherent calibrated port tuple, the constructed polynomial has a
zero in ten strictly positive integers if and only if there is a
finite syntactic ZFC refutation.  Consequently ZFC proves that its
lack of positive integer solutions is equivalent to
$\operatorname{Con}(\mathrm{ZFC})$.
\end{theorem}

\begin{proof}
Soundness and completeness were proved above, using the same
arithmetizable source calibration as the parent.  Negating the
resulting existence equivalence gives the asserted consistency
equivalence.  The argument is ordinary mathematical proof
formalizable in ZFC; it is not an assertion that a new Lean, Rocq,
or encoded internal-ZFC proof artifact has been produced.
\end{proof}


\section{Exact degrees and highest homogeneous forms}
\label{sec:saturated-two-root}

This section proves the global degree bounds and their attainment for
the four literal circuits.  Write $\gamma=\varepsilon$ for the fixed
center choice; it is not another unknown.  The square-recovery argument
and the complete witness constructions were proved in the preceding
section.

\subsection{Literal degrees and highest homogeneous forms}

Set $q_0=\beta J$ and $k_0=\eta+\zeta$.  With the port values fixed
as constants, the four exact results are
\[
\begin{array}{c|r|r}
\text{circuit}&\text{live binary operations}&\text{total degree}\\\hline
 X=wq,\ C_1=(A+1)UE&92&283\\
 X=wq,\ C_0=AUE&93&278\\
 X=q+x,\ C_1=(A+1)UE&98&258\\
 X=q+x,\ C_0=AUE&99&254.
\end{array}
\]
The operation counts refer to the supplied literal graphs after
discarding unused nodes; they charge additions, subtractions and
multiplications once each, including multiplication by a fixed integer.

\paragraph{Exact-degree certificate.}
For the $X=wq$ circuits, the highest homogeneous forms are
\begin{align*}
 [N_{2,1}]_{283}
   &=2^{11}q_0^{149}w^{29}Y^{61}i^{10}k_0^{32}(q_0-F)^2,\\
 [N_{2,0}]_{278}
   &=2^{10}q_0^{146}w^{28}Y^{60}i^{10}k_0^{32}(q_0-F)^2.
\end{align*}
For the direct circuits, put $x_0=q_0+x$.  The highest forms are
\begin{align*}
 [N_{2,1}^{\rm dir}]_{258}
   &=2^{11}q_0^{124}x_0^{29}Y^{61}i^{10}k_0^{32}(q_0-F)^2,\\
 [N_{2,0}^{\rm dir}]_{254}
   &=2^{10}q_0^{122}x_0^{28}Y^{60}i^{10}k_0^{32}(q_0-F)^2.
\end{align*}
These are nonzero polynomials for every fixed $\beta>0$.  Therefore
they certify exact degrees for the actual coefficient tuple.

\paragraph{Proof.}
Scale all ten witnesses by $t$, keep the coefficient ports fixed, and
work in the formal Laurent-series field at $t=\infty$ over the rational
function field in the ten unscaled witnesses.  A series degree is the
largest exponent of $t$.  For polynomials it is their ordinary total
degree.  Choosing formal square roots with the indicated square leading
forms is justified by the binomial series; each relative correction
has negative degree.  No analytic convergence is used.

In either source the literal numerator satisfies
\[
 Z:=Z_{\rm num}/M=V+C_\gamma+\mathcal E,
\]
where the rational remainder has degree at most $4$.
For $X=wq$, it is
\[
 \mathcal E=n-LB_t+K_{\rm pos}/H+Y/q^3-B_t/(q-1)-e/\mu^2,
 \quad L=k(T+\alpha),\quad H=4A-5,
\]
and $M=q^3(q-1)\mu^2H$.
For $X=q+x$, replace the remainder by
\[
 \mathcal E=n-LB+K_{\rm pos}/H+Y/q^3
                -B/[q(q-1)]-X/q^4-e/\mu^2,
\]
and retain the extra factor of $q$ in
$M=q^4(q-1)\mu^2H$.  Both denominators have degree $13$.
These rational identities are exact readings of the respective
polynomial graphs, not additional defining equations.

Choose formal roots $g^2=ED_f$ near $AUE$ and
$y^2=D_y/\delta$ near $V$.  The two exact identities are
\[
 (AUE)^2-ED_f=U^2E^2-E,                         \tag{D1}
\]
and
\[
 (y-V)(y+V)=(V^2-1)/\delta.                       \tag{D2}
\]
The degrees controlling the argument are
\[
\begin{array}{c|rrrrrrrrrr}
 &A&U&E&D_f&\delta&V&M&g&C_1-g&C_0-g\\\hline
 X=wq&5&5&18&38&48&42&13&28&23&18\\
 X=q+x&4&5&16&34&42&38&13&25&21&17.
\end{array}
\]
Equation (D2) gives $\deg(y-V)=-6$ and $-4$, respectively.
Equation (D1) gives
\[
 C_0-g=\frac{U^2E^2-E}{AUE+g},
 \qquad C_1-g=UE+(C_0-g).
\]
The leading forms of these two defects are $u_0e_0/(2a_0)$ and
$u_0e_0$, respectively, where the subscripts denote leading forms.
In every case the defect degree is greater than $4$, so it dominates
both $\mathcal E$ and $V-y$.

The exact four-factor norm identity is
\[
 N_2=\delta^2M^4
       \prod_{\sigma,\rho\in\{-1,1\}}(Z+\sigma y+\rho g).
                                                               \tag{D3}
\]
The two factors with $\sigma=1$ have leading form $2v_0$.
The factor with $\sigma=-1,\rho=1$ has leading form $2g_0$.
The last factor has leading form $u_0e_0$ for $C_1$ and
$u_0e_0/(2a_0)$ for $C_0$.
Thus the four degree totals are exactly
\begin{align*}
 2(48)+4(13)+2(42)+28+23&=283,\\
 2(48)+4(13)+2(42)+28+18&=278,\\
 2(42)+4(13)+2(38)+25+21&=258,\\
 2(42)+4(13)+2(38)+25+17&=254.
\end{align*}

To verify the coefficients without any large-numeral expansion, for
$X=wq$ put
\[
 a_0=q_0^3wY,\quad u_0=ik_0^2Y^2,\quad
 e_0=q_0^{10}w^2Y^4k_0^2,\quad r_0=q_0^3(q_0-F).
\]
Then
\[
 d_0=a_0^4u_0^2e_0,\quad m_0=4q_0^8a_0,\quad
 v_0=-a_0^2u_0^2e_0r_0,\quad g_0=a_0u_0e_0.
\]
Multiplying the leading factors in (D3) gives
\[
 8d_0^2m_0^4v_0^2g_0u_0e_0
 \quad\hbox{or}\quad
 4d_0^2m_0^4v_0^2u_0^2e_0^2,
\]
respectively.  These simplify to the two asserted $X=wq$ forms.
For the direct source replace
\[
 a_0=q_0^2x_0Y,\qquad
 e_0=q_0^8x_0^2Y^4k_0^2,\qquad m_0=4q_0^9a_0;
\]
the remaining abbreviations retain their formulas.  The same products
give the two direct forms.  Every displayed final top form is a nonzero
polynomial, proving attainment as well as the upper bound.
\hfill$\square$

\paragraph{Independent exact polynomial replay.}
The standard-library verifier evaluates every gate of each supplied
graph as a univariate polynomial over the integers.  On the ray where
all ten witnesses equal $t$ and
$(\beta,K_0,M_C,M_F)=(2,3,5,7)$, it obtains, in table order,
\[
 2^{192}t^{283},\qquad 2^{188}t^{278},\qquad
 2^{167}3^{29}t^{258},\qquad 2^{164}3^{28}t^{254}
\]
as the exact leading terms.  This is finite arithmetic verification of
one specialization.  The global Laurent-series proof, including its
explicit nonzero top homogeneous forms, is what establishes the claims
for the real fixed coefficient tuple.


\section{A saturated equation over twelve unrestricted integers}
\label{sec:integer12-saturated}

The polynomial in \path{integer12_saturated/rational.json} has twelve
unrestricted integer unknowns,
123 binary arithmetic operations, and exact total degree 1084.  It
strictly improves both operation count and degree over the pinned
scalar-carrier parents, which have metrics $145/1200$ and $141/1456$
\cite{scalar12}.
The construction combines coprime Pell saturation with the scalar
intersective carrier.  Its completeness requires a rescaling of the
auxiliary coordinate and the full arbitrary-residue interpolation theorem.
The new rescaling and its compatibility with that retained theorem
are proved below.

\subsection{Complete twelve-integer polynomial}

The independent coordinates are
\[
 (j,C,z,h,r,v,w,T,c,n,t,s)\in\Z^{12}.
\]
Use the exact source34 \texttt{mixed-parity-unrounded-1064921} coefficient
tuple, separately from the compact44 tuple of the positive-domain forms:
\[
 \beta=2^{1064921}-1,\qquad
 \kappa_1=2^{1113},\qquad \kappa_2=2^{666104}.
\]
The remaining fixed integers $K_0,M_C,M_F$ are fully specified by the
included six payloads and the finite source34 recipe, complete tile table,
and exact machine seed.  The supplied $M_F$ is already shifted.  Every
coefficient is reconstructed and compared byte for byte by the checker.
No free coefficient parameter remains.

Define the following ordinary polynomial abbreviations; none is another
unknown.  In this section $E$ denotes the first Pell radicand.
\begin{align*}
 J&=6j+1,&q&=\beta J+1,&W&=12w+4,\\
 X&=qW,&Y&=(2v+1)q^3,&a&=(X+1)Y,\\
 H&=4a+3,&\Delta&=a^2+H,&F&=(K_0+W)C-z(q-1).
\end{align*}
\begin{align*}
 R&=\bigl(q(q-F)-C\bigr)(q^2-1)+(M_C+qM_F)J,\\
 k&=R+1+hXY,\qquad \eta=c-Yk,\qquad\zeta=k-\eta,\\
 G&=(\kappa_1F)(\kappa_2C)-(q-\kappa_1F-\kappa_2C)^2,\\
 \Pi&=G\eta\zeta-(t^2-G-\eta\zeta+1)^2,\qquad
 D_g=(3\Pi-4)s^2+1,\\
 \ell&=XY^2,\qquad E=1+\ell(\ell+1)k^2,\qquad U=src^2,\\
 D_f&=1+\Delta U^2E,\qquad Q=\Delta^2U^2E,\qquad\delta=Q-1,\\
 b&=c-R+4UET,\qquad V=D_fb-c,\qquad C_y=QV^2-1,\\
 K&=X(X+2ac)-1,\qquad S=UH^5,\qquad z_c=HT+K,\\
 A_c&=(z_c^2+3H^2)(z_c^3-19H^3),\qquad Z=S(n+V)-A_c.
\end{align*}
Put
\begin{align*}
 f_s&=\delta Z^2,& f_a&=9\delta S^2ED_f,\\
 f_b&=81\delta S^2D_g,& f_c&=S^2C_y,\\
 P_e&=(f_s+f_c-f_a-f_b)^2-4(f_sf_c+f_af_b).
\end{align*}
The entire requested equation over $\Z^{12}$ is
\begin{equation}
 \boxed{\quad P=P_e^2-64f_sf_cf_af_b=0.\quad}
 \label{eq:saturated-integer12}
\end{equation}
Every operation is polynomial arithmetic.  The literal circuit evaluates
$81D_g$ directly as $(243\Pi-324)s^2+81$, saving one multiplication.
Fixed powers in the display are implemented by counted binary
multiplications.

For interpretation, let $\mathcal F_8(Z;A,B,C)$ be the product of
$Z-\epsilon_1\sqrt A-\epsilon_2\sqrt B-\epsilon_3\sqrt C$ over the eight
sign choices.  Sign invariance makes this an integer polynomial.  Then
\begin{equation}
 P=\delta^4\mathcal F_8\left(
 Z;9S^2ED_f,81S^2D_g,\frac{S^2C_y}{\delta}\right).
 \label{eq:saturated-integer12-norm}
\end{equation}
This is only a norm interpretation of the division-free
definition~\eqref{eq:saturated-integer12}.  The denominator is cleared
symbolically, and $\delta$ will be shown never to vanish on integer inputs.

\subsection{Signed bootstrap, separation, and source restoration}

\begin{theorem}[Twelve-integer saturated carrier]
Under the same ordinary native, compiler, and refutation calibration as the
pinned scalar-carrier parent,
\[
 \ZFC\vdash\left(
   \forall x\in\Z^{12}\;P(x)\ne0
 \right)\longleftrightarrow\Con(\ZFC).
\]
The new arithmetic equivalence is proved below and is formalizable in
ZFC.  This statement does not claim a new encoded object-ZFC derivation
or proof-assistant endpoint for the entire inherited calibration.
\end{theorem}

\begin{proof}[Soundness]
The fixed parity and residue choices give, unconditionally,
\[
 |q|\ge16,\quad |X|\ge64,\quad |Y|\ge4096,\quad
 H\ne0,\quad\Delta=(a+2)^2-1>1.
\]
Also $E\ge1$, since $\ell(\ell+1)\ge0$ for integer $\ell$.
If $U=0$, then $Q=0$; otherwise $Q\ge\Delta^2>1$.  Thus $\delta\ne0$
for every integer assignment.

Write $\pi(x)=(x^2+3)(x^3-19)$.  If $U=0$, then $S=0$ and
the final polynomial becomes $A_c^8$.  But
$A_c=H^5\pi(T+K/H)\ne0$: the quadratic has no real zero and the cubic
has no rational zero.  Hence a zero of $P$ has $U\ne0$, so $s,r,c$
are nonzero and $|U|\ge c^2$.  We obtain
$D_f>|c|+1$.  Since $b$ is integral, either $b=0$ and $|V|=|c|\ge1$,
or $|V|\ge D_f-|c|>1$.  Consequently $C_y>0$ and $Q>1$ before any
guard or projection is assumed.

Divide a vanishing conjugate factor of~\eqref{eq:saturated-integer12-norm}
by $S$.  If $D_g<0$, its summand is the only nonzero imaginary term,
an impossibility.  Thus $D_g\ge0$.  Since $s\ne0$, the inequality
$\Pi\le0$ would force $D_g\le-3$; hence $\Pi>0$.  This implies both
$G>0$ and $\eta\zeta>0$.  Indeed their product is positive; if they
were both negative, their absolute values $A,B>0$ would give
$\Pi=AB-(t^2+A+B+1)^2<0$.

Modulo 3, $X=-1$ and $a=H=\Delta=0$.  The value $\ell=XY^2$ is
0 or $-1$, so $E=D_f=1$ and $\delta=C_y=-1$ modulo 3.  The three
normalized roots have squared 3-adic orders
\[
 \sqrt{C_y/\delta}:0,\qquad
 3\sqrt{ED_f}:2,\qquad
 9\sqrt{D_g}:4+v_3(D_g).
\]
Omit a zero guard term.  In a rational sum of square roots of positive
rationals, each nontrivial square class must cancel internally.  Within
such a class the ratios are rational, and the displayed valuations
give a unique term of minimum 3-adic order, making cancellation
impossible.  Therefore all three roots are rational.

The integer $ED_f$ is thus a square.  Since
$D_f=1+\Delta U^2E$, we have $\gcd(E,D_f)=1$; positivity and unique
factorization give
\[
 E=\tau^2,\qquad D_f=f^2,\qquad I=U\tau,
 \qquad \tau,f\in\Z_{>0}.
\]
This includes $E=1$.  We have $c^2\mid I$, $I\ne0$,
$f^2=1+\Delta I^2$, and $Q=(\Delta I)^2$.
The rational auxiliary root $y$ satisfies
\begin{equation}
 QV^2-(Q-1)y^2=1.
 \label{eq:sat12-aux}
\end{equation}
The adjacent-coefficient descent from the retained parent makes $y$
integral.  Explicitly, if $|y|=a/b$ in lowest terms then $b^2\mid Q-1$;
the norm-preserving map
\[
 V'=(2Q-1)V-2(Q-1)y,\qquad y'=(2Q-1)y-2QV
\]
preserves the reduced denominator and strictly decreases $|V|$ until
$|V|=1$, where $|y|=1$.  Signs are reflected during the descent.
The guard root is integral as well.

The vanishing factor now gives
\[
 S(n+V)-A_c=S(\pm3\tau f\pm9g\pm y),
\]
so $S\mid A_c$ and $H^5\mid A_c$.  Expanding
\[
 \Phi_H(z)=z^5+3H^2z^3-19H^3z^2-57H^5,
\]
if $p^e\Vert H$ and $v_p(z)=b<e$, its leading term has valuation
$5b$ and every other term has larger valuation.  Thus
$H^5\mid\Phi_H(z)$ implies $H\mid z$.  Applying this to $z=HT+K$
recovers $H\mid K$, with no sign assumption on $H$.

The exact source coordinates are
\begin{equation}
 r_{\rm old}=r\tau,\qquad T_{\rm old}=\tau T,
 \qquad I=sr_{\rm old}c^2.
 \label{eq:sat12-restoration}
\end{equation}
They give
\[
 4IT_{\rm old}=4UE T,\qquad
 V=f^2(c-R+4IT_{\rm old})-c.
\]
All strict guards, the first and strong norm squares,
the auxiliary norm~\eqref{eq:sat12-aux}, $c^2\mid I$, and $H\mid K$
are now exactly the hypotheses of the pinned signed scalar-carrier
restoration theorem.  That theorem permits arbitrary nonzero strong
multipliers, including even ones.

For precision, it uses absolute Pell folding and Lucas multiplication
to recover $|c|=\psi_{|a+2|}(R)$; the auxiliary residues and strict
index bounds identify the index.  Its whole-product estimate, using
$H\mid K$, then yields $\operatorname{sign}(c)X=2^R$ and restores
the missing main root and generalized signed native kernel.  These
unchanged lemmas are retained in \path{integer12_saturated/parent/PROOF.md}
and its pinned signed proof sources.  The named native/refutation
calibration therefore supplies a finite syntactic ZFC refutation.
The claim is restoration of that kernel, not a pointwise substitution
in the old scalar polynomial, whose carrier uses $T_{\rm old}$.
\end{proof}

\subsection{Completeness with a fixed strong witness}

Start with the canonical native witness of a finite ZFC refutation.
Its relevant conditions are
\[
 A_0=a+2>1\text{ even},\quad R\equiv3\pmod4,\quad
 c=\psi_{A_0}(R)>0,\quad c\equiv R\pmod4,
 \quad E=\tau^2,\quad\tau>0.
\]
Choose $t$ for the joint strict guard, then a positive $s$ with
$D_g=g^2$, as in the parent.  When $\Pi=1$ one can use $s=1,g=0$;
otherwise ordinary Pell existence supplies the guard witness.

The full-residue interpolation theorem in the retained
\path{integer12_saturated/parent/source/INTERPOLATION12.md}, Section 6, says that for
\emph{every} prescribed positive $N$, there is one fixed strong pair
$I,f$ with $4Nc^2\mid I$, such that \emph{every} residue class modulo
that same $I$ is attained by an auxiliary $T_{\rm old}$ satisfying the
norm.  Its strong pair is chosen before the residue request, and its
squared-modulus period is $4mf$, where $m$ is the strong Pell index.

Apply this theorem with $N=s\tau$, obtaining
\[
 4s\tau c^2\mid I,\qquad U=I/\tau,
 \qquad r=I/(s\tau c^2).
\]
Thus $U=src^2$ is integral.  The canonical main projection gives
$H\mid K$; put $\kappa=K/H$.  The retained intersectivity theorem for
$\pi(x)=(x^2+3)(x^3-19)$ provides a residue $a_*$ modulo $U$ with
\[
 \pi(a_*+\kappa)\equiv0\pmod U.
\]
Now request the old auxiliary residue
\[
 T_{\rm old}\equiv\tau a_*\pmod I.
\]
Since $I=\tau U$, this ensures $\tau\mid T_{\rm old}$ and, on setting
$T=T_{\rm old}/\tau$, gives exactly $T\equiv a_*\pmod U$.
The new auxiliary expression is therefore the required source
expression, and
\[
 \frac{A_c}{S}=\frac{\pi(T+\kappa)}{U}\in\Z.
\]
Finally choose
\[
 n=\frac{A_c}{S}+3\tau f+9g+y-V\in\Z.
\]
This makes a conjugate factor vanish and proves completeness.
No positivity of $n$, coprimality of $\tau$ and $U$, or change of strong
pair after the residue choice is needed.  The factor $E$ in $4UET$ is
essential to this exact quotient-coordinate argument.

\subsection{Attained degree and reproducible cost}

All fixed coefficients have degree zero.  The relevant total degrees are
\[
\begin{array}{c|rrrrrrrrrr}
 &U&E&\Delta&D_f&\delta&b&V&C_y&S&A_c\\\hline
 \deg&4&34&12&54&66&39&93&252&34&45.
\end{array}
\]
Put $z_0=Sn-A_c$, of degree 45, and $\omega=f_s-f_c$.
The defining relation $C_y=(\delta+1)V^2-1$ gives the exact cancellation
\[
 \omega=\delta z_0^2+2\delta SVz_0-S^2(V^2-1).
\]
Its unique top term is $-S_0^2V_0^2$, of degree 254; competing degrees
are 156, 238, and 68.  The four squared arguments have degrees
\[
 \deg f_a=222,\qquad\deg f_b=180,\qquad
 \deg f_c=\deg f_s=320.
\]
The identity
\[
 P_e=\omega^2-2(f_a+f_b)(2f_c+\omega)+(f_a-f_b)^2
\]
has unique top term $-4(f_a)_0(f_c)_0$, of degree 542; every competitor
has degree at most 508.  The correction $64f_sf_cf_af_b$ has degree
1042, so the final degree is $2\cdot542=1084$, provided the top form
does not vanish at the fixed tuple.

It does not.  With subscript zero denoting highest homogeneous part,
\begin{align*}
 q_0&=6\beta j,&X_0&=12q_0w,&Y_0&=2q_0^3v,\\
 a_0&=X_0Y_0,&U_0&=src^2,&E_0&=h^2X_0^4Y_0^6,\\
 S_0&=2^{10}a_0^5U_0,&\delta_0&=a_0^4U_0^2E_0,
   &(D_f)_0&=a_0^2U_0^2E_0,\\
 V_0&=4a_0^2U_0^3E_0^2T.&&&&
\end{align*}
The term $4UET$ uniquely dominates $R$ in $b$, so these are single
monomials.  The entire degree-1084 homogeneous part is
\begin{equation}
 P_0=2^{1144}3^{760}\beta^{624}
 j^{624}w^{132}v^{164}h^{32}s^{32}r^{32}c^{64}T^4\ne0.
 \label{eq:sat12-top}
\end{equation}
The exponents sum to 1084 and the actual fixed $\beta$ is positive.
Thus the degree is proved for the complete polynomial, independently
of finite line tests.

The literal operation count is 75 multiplications, 25 additions, and
23 subtractions: 123 in total, including scalar multiplications.  Every
arithmetic node is live.  The 13-operation final norm follows a
110-operation prefix and carrier.  An included integral control has the
same operation count and degree 1280; it is dominated by this candidate.

The independent standard-library checker imports no producer.  It
verifies every prefix by free-ring identities, all 35 coefficients of
the eight-conjugate norm, the centered cancellation and restoration
identities, exact leading induction, complete coefficient bytes and
recipe, and operation counts.  Normal and optimized runs pass; corrupted
auxiliary, strong, guard, carrier, and final expressions are rejected.
These computations certify the literal polynomial and its algebraic
metrics.  The signed semantic argument and interpolation proof are
ordinary mathematical proofs using the explicitly retained source
interfaces, independently reviewed in the accompanying audit.


\section{Proof dependencies, reproducibility, and limitations}
\label{sec:dependencies}

\subsection{The inherited semantic boundary}

The new positive proofs reuse the following ordinary theorems from the
pinned construction.  Their exact statements and source proofs are
retained under \texttt{source/proofs/}:
\begin{enumerate}
\item Main/auxiliary Pell folding recovers $c=\psi_A(p)$ from the strong
      equation, the auxiliary equation, and their retained congruences.
\item The strong rank and signed auxiliary-index lemmas recover $p=R$
      once the transport argument has supplied the correct rank range.
\item The zero-offset Pell-ratio estimate, whole-product projection,
      fixed-minus parity, and reconstruction of the unscaled source
      apply to the recovered outer data.  In the direct variant, delayed
      $q$-divisibility is part of this interface.
\item The complete native/controller theorem, together with the K215
      syntactic-refutation calibration and the coherent BFS state/color
      transfer, identifies source solutions with finite ZFC refutations.
\end{enumerate}
The fourth item is a substantive inherited theorem.  Reconstructing the
coefficient bytes or running a finite machine prefix would not prove it.
This packet does not describe those finite computations as such a proof.
The new arithmetic arguments are uniform for the specified port box;
the final fixed tuple is the \emph{whole} coherent BFS tuple from the
complete source page.  Its $K_0$ is not equated with the different original
compact44 $K_0$.  No individual port is substituted between tuples.

The twelve-integer packets have their own six-port source34 tuple,
\texttt{mixed-parity-unrounded-1064921}.  The parent and dependencies are
retained separately under \texttt{integer12\_saturated/parent/} and
\texttt{integer12/parent/}.  The two coefficient
families are not mixed.  All ordinary proofs use only finite algebra,
integer arithmetic, induction, and the previously arithmetized source
calibration.  The remaining end-to-end proof-assistant work is explicitly
outside this contribution.

\subsection{What the executable checks establish}

\begin{itemize}
\item Every literal gate, input, fixed port, and live operation count is
      checked; deterministic rebuilds reproduce the exported JSON files.
\item The complete generic two-root norm has six terms after expansion;
      its equality with the four signed factors is verified exactly.
      The intermediate and saturated integer three-root norms, and the
      earlier twelve-integer sixteen-factor norm, are checked independently
      as well.
\item The anti-cancellation congruence and shared-product identities are
      checked over free polynomial rings, independently of sample values.
\item The complete four compact44 coefficients are reconstructed from
      the printed 436-character seed and all 943 tiles, then compared
      against every supplied coefficient byte.  The twelve-integer
      coefficient recipe is replayed separately.
\item Exact univariate polynomial replays corroborate all four saturated
      degree certificates.  The global upper bounds and nonzero highest
      homogeneous forms are proved symbolically in the article; they
      are not inferred from those one-dimensional checks.
\item Independently transcribed complete formulas are bound to every
      gate of the four positive saturated graphs and the twelve-integer
      saturated graphs by exact identities over free local polynomial
      rings.  Modular evaluations with the actual fixed coefficients
      provide additional transcription controls.  These checks do not
      prove integer unsolvability.
\end{itemize}

From the extracted archive root, the main finite replay commands are
\begin{verbatim}
python3 -B build.py --check
python3 -B verify.py
python3 -B -O verify.py
python3 -B integer12_saturated/build.py --check
python3 -B integer12_saturated/verify.py
python3 -B -O integer12_saturated/verify.py
python3 -B integer12/build.py --check
python3 -B integer12/verify.py --output /tmp/integer12-check.json
python3 -B alternative/relational_kernel.py
latexmk -pdf -interaction=nonstopmode -halt-on-error report.tex
\end{verbatim}
Only the Python standard library is needed for the arithmetic checks.
LaTeX is needed only to rebuild the PDF.  No network access is needed
after extraction.  Verification receipts and provenance hashes are
included, together with the complete sources.

\section{Further research questions}

\begin{enumerate}
\item \textbf{Can saturation be applied to more square constraints without
      increasing witness-reconstruction cost?}  Coprimality alone is not
      enough: every absorbed factor must be supplied by a compatible
      canonical witness.  The Pell-period argument here solves that
      problem for the first and strong ordinates.
\item \textbf{Can the upper center be approximated more accurately by a
      cheap integer polynomial?}  The sharp center $AUE$ leaves defect
      $(U^2E^2-E)/(AUE+\sqrt{ED_f})$.  Removing its highest homogeneous
      part could lower degree further, but extra denominators and their
      exact integer recovery must be counted.
\item \textbf{Can a nine-positive form use the same two-root interface?}
      Eliminating an additional interpolation coordinate requires a
      proved simultaneous Pell-divisibility construction.  Setting a
      witness to $1$, or sharing two independent ordinates without that
      theorem, does not establish completeness.
\item \textbf{Can a source based directly on equational proofs beat the
      native history compiler?}  The local kernel in the final appendix
      is explicit and small.  A fully printed finite axiom tuple, a
      conservative calibration, and a counted unbounded-history compiler
      are still needed before it becomes a competing consistency equation.
\item \textbf{Which ordinary calibration links should be formalized next?}
      A useful next target is the new product-square recovery and the
      finite Pell-period completeness map, followed by composition with
      the existing native/internal-arithmetic work.  Rechecking only
      finite evaluations would not discharge that task.
\end{enumerate}

\appendix
\section{The complete fixed coefficient recipe}\label{sec:constants}
The following finite recipe is copied from the complete, calibrated BFS44 page at the pinned snapshot. Every sum, index, floor, alphabet symbol, and tile in this appendix defines a fixed integer before the polynomial is evaluated. The independent reconstruction in \texttt{coefficients/recipe\_compact44.py} reproduces every coefficient byte.
\subsection*{Four integers}
Every index, set, sum, product and floor below defines a constant, not
another unknown. Put
\[
\rho=2^{67},\quad B=\rho^{8807},\quad\beta=B-1,\quad
G_m(x)=\sum_{j=0}^{m-1}x^j.
\]
Define
\[
\begin{split}
\mathcal E={}&\rg0{942}\cup\rg{983}{992}\cup\rg{1003}{1011}\\
&\cup\rg{1013}{1031}\cup\{1063,1064,1095,2195,3296\}.
\end{split}
\]
Then
\begin{align*}
M_C={}&B-1-\sum_{e\in\mathcal E\setminus\{1\}}\rho^e-2\rho^{1063},\\
M_F={}&B-1+4+(2^{51}-2)(\rho^{4393}+\rho^{5495})\\
&+\rho^{1013}G_{10}(\rho)+\rho^{1065}G_9(\rho)
  +\rho^{2195}+\rho^{3296}\\
&+(2^{65}-1)\rho^{8792}G_{13}(\rho)+(2^{19}-1)\rho^{8805},\\
K={}&\rho^{30}+\rho^{62}+\rho^{1100}+\rho^{2201}+\rho^{8849}\\
&+(\rho^{4393}+\rho^{5495})\Phi(1/\rho),
\end{align*}
where
\[
\begin{split}
\Phi(x)={}&(1+2^{51})G_{943}(x)
 -2^{10}\sum_{t\in\mathcal T}v(t)x^{\varrho(t)}\\
&+\sum_{(b,e,m)\in\mathcal C}2^b x^e G_m(2x),\\
\mathcal C={}&\{(10,1023,8),(18,1013,9),(27,1003,8),\\
&(35,983,9),(44,1031,1),(45,1022,1),(46,1011,1),\\
&(47,992,1),(48,1095,1),(49,2195,1),(50,3296,1)\}.
\end{split}
\]
The rational intermediate in K cancels exactly; all four results are positive integers.

\subsection*{The 943 tiles and their ranks}
Use the actions $u_j,v_j,d_j$ decoded below. Let
\[
\begin{gathered}
h(q,b)=2(q+1)+b,\quad
\mathcal R=\{u_j:d_j=1,\ 0\le j<430\},\\
\mathcal L=\{u_j:d_j=0,\ 0\le j<430\},\\
r_q=6+|\mathcal R\cap\rg0{q-1}|,\quad
l_q=205+|\mathcal L\cap\rg0{q-1}|.
\end{gathered}
\]
An interval ending at $-1$ is empty. The following union is $\mathcal T$;
the index ranges in each row are fixed:
\[
\setlength{\arraycolsep}{1.4pt}\begin{array}{c|l}
\text{range}&\text{tile }(w,n,e,s)\\\hline
j=2q+b<430&((1-d_j)l_{u_j},h(q,b),d_jr_{u_j},v_j)\\
b=0,1&(0,h(215,b),0,h(215,b)),\ (0,b,0,b)\\
q\in\mathcal R,\ b=0,1&(r_q,b,0,h(q,b))\\
q\in\mathcal L,\ b=0,1&(0,b,l_q,h(q,b))\\
t\in\{0,1,432,433\}&(0,t,0,438)\\
&(2,434,2,0),\ (2,435,3,1)\\
&(3,436,4,h(u_0,0)),\ (4,436,4,0),\ (5,0,0,0)\\
t=434,435,436&(1,438,1,t)\\
a\in\{0,1,4\},\ b=0,1,2&(a,437,b,437)
\end{array}
\]
Put $s_*=(2,435,3,1)$, $f_*=(5,0,0,0)$, and $\mathcal T'=\mathcal T\setminus\{s_*,f_*\}$.
For $t=(w,n,e,s)$ set
\[
\begin{gathered}
\kappa(t)=2^{30}w+2^{20}n+2^{10}e+s,\\
v(t)=w+2^8n+2^{17}e+2^{25}s+127\,2^{34},\\
\varrho(s_*)=0,\quad\varrho(f_*)=1,\\
\varrho(t)=2+\sum_{z\in\mathcal T'}
 \fl{\frac{2^{40}+\kappa(t)-\kappa(z)-1}{2^{40}}}
 \quad(t\in\mathcal T').
\end{gathered}
\]

\subsection*{A finite decoder for the 430 actions}
Put $D_b(A,k)=\fl{A/b^k}\bmod b$, where
$a\bmod b=a-b\fl{a/b}$. For $0\le q<215$, $b=0,1$, $j=2q+b$:
\[
\begin{gathered}
N=\mathcal A\bmod2^{636},\quad
T_0=\fl{\mathcal A/2^{636}}\bmod2^{1675},\\
E_0=\fl{\mathcal A/2^{2311}},\\
\nu_j=D_2(N,j),\quad \sigma_j=\sum_{k=0}^{j-1}\nu_k,\\
m_q=D_2(N,430+q),\quad \tau_q=\sum_{k=0}^{q-1}m_k,\\
F_q=14(1-m_q)+m_qD_{16}(E_0,\tau_q),\\
u_j=\nu_j(1+\sigma_j)+(1-\nu_j)D_{216}(T_0,j-\sigma_j),\\
v_j=D_2(F_q,2b),\quad d_j=D_2(F_q,2b+1).
\end{gathered}
\]
Empty sums are zero. Here $(v_0,d_0,u_0)=(1,1,1)$; the constructor tile
uses this actual first target, including its state renaming.

\subsection*{The integer $\mathcal A$}
The alphabet, in order of digit values $0,\ldots,61$, is
\begin{center}\small\ttfamily
0123456789ABCDEFGHIJKLMNOPQRSTUVWXYZabcdefghijklmnopqrstuvwxyz
\end{center}
Concatenate the following 436 characters as digits $e_0,\ldots,e_{435}$;
put $\mathcal A=\sum_{j=0}^{435}e_j62^{435-j}$.\par
\begingroup\small\ttfamily
BxEnhvy6NipzkWmROo0GsA3uK3twWLxscLq20UTumy5Cdu461LwD6g7RDNzF8Y\par
iYDFTpB8ehAypd7d6OsOL4APamnJeNXQypUezCmW9n4oXP6XfgmiWuAWZ8ZJ7U\par
OP7gu4kVdlGZqZHULQMar1dXNeJW4i7oNVI8usp8BlovjjZbwPCovZI328fRUq\par
4q92R81s5GjWQaQfGhkq9UYMK3rRMaBmsOIBkJqc2cj2VRIZ6mjzkBIOkW9SjZ\par
FzeA6J5u4j3sBL3SdgwXzM9mHvQivu8EJ1TBTpRhHDqZIbmYgwQ8VqTAFotl6F\par
We1YC554GKrVHsA8oUGKSA6F6k5KREBSD0Tgd7bFALY84LLbpw574BQyvdoHrl\par
0Oe35uwDf15DpT6VJWZBjSbHPpPAINYTte0eUxnTFTKBlGj8Am6sb1gLldh8WZ\par
8P\par
\endgroup

The archive uses little-endian unsigned coefficient bytes, compressed with gzip. The bit lengths of $\beta,M_C,M_F,K_0$ are respectively $590069,590069,590070,592884$. Here the appendix symbol $K$ is the polynomial port $K_0$; $M_F$ already includes its prescribed shift by $\beta$. The accompanying payloads constitute one coherent tuple.

\section{Coordinate-shear controls over twelve integers}
\label{sec:integer12-main}

The following two earlier improvements are retained as simple algebraic
controls.  The saturated equation in Section~\ref{sec:integer12-saturated}
has a better operation count and degree than either.  These controls
use a different proof mechanism: invertible polynomial coordinate
changes covering the entire integer domain, including exceptional
values of the denominators used only in explanations.

\label{sec:integer12}

The preceding positive-domain constructions and the following
integer-domain constructions are separate points in the comparison.
Here all twelve unknowns range over $\Z$.  We give two polynomial
coordinate changes of the pinned repository's current scalar-carrier
equations.  The transformations preserve the entire integer zero sets
bijectively and improve their degree--cost tradeoffs.  Centering is already
used elsewhere in the repository; the contribution here is its application
to the present scalar four-root carrier, with complete circuits and attained
degree certificates.

\subsection{Fixed coefficients and complete equations}

The ordered unknowns are
\[
 (j,C,z,h,r,v,w,T,c,n,t,s)\in\Z^{12}.
\]
Use the coherent \texttt{mixed-parity-unrounded-1064921} tuple, independently
of the compact44 tuple used by the positive-domain constructions:
\[
 \beta=2^{1064921}-1,\qquad
 \kappa_1=2^{1113},\qquad \kappa_2=2^{666104}.
\]
The remaining fixed integers $K_0,M_C,M_F$ are the exact source34 values
in the included coefficient payloads.  The complete finite recipe,
literal machine seed, and all 943 tiles are included; the verifier rebuilds
every coefficient and compares all bytes.  In particular $M_F$ is already
shifted by $\beta$.  The six coefficients are never mixed with another
tuple.

Define the following polynomial abbreviations; none is an additional
unknown:
\begin{align*}
 J&=6j+1,&q&=\beta J+1,&W&=12w+4,\\
 X&=qW,&Y&=(2v+1)q^3,&E&=XY,\\
 a&=E+Y,&H&=4a+3,&\Delta&=a^2+H,\\
 F&=(K_0+W)C-z(q-1),&&&
\end{align*}
\begin{align*}
 R&=\bigl(q(q-F)-C\bigr)(q^2-1)+(M_C+qM_F)J,\\
 k&=R+1+hE,\qquad \eta=c-Yk,\qquad \zeta=k-\eta,\\
 G&=(\kappa_1F)(\kappa_2C)-(q-\kappa_1F-\kappa_2C)^2,\\
 \Pi&=G\eta\zeta-(t^2-G-\eta\zeta+1)^2,\\
 D_g&=(3\Pi-4)s^2+1,\qquad I=src^2,\\
 D_f&=1+\Delta I^2,\qquad Q=\Delta^2I^2,\qquad\delta=Q-1,\\
 V&=D_f(c-R+4IT)-c,\qquad C_y=QV^2-1,\\
 \ell&=XY^2,\qquad D_t=1+\ell(\ell+1)k^2,\\
 K&=X(X+2ac)-1,\qquad S=IH^5,\qquad z_c=HT+K,\\
 A_c&=(z_c^2+3H^2)(z_c^3-19H^3).
\end{align*}
Set
\[
 r_a=9S^2D_f,\qquad r_b=81S^2D_t,\qquad r_e=729S^2D_g.
\]
Let $\mathcal F_{16}(Z;u_1,u_2,u_3,u_4)$ be the product of
$Z-\sum_{i=1}^4\epsilon_i\sqrt{u_i}$ over the sixteen choices
$\epsilon\in\{-1,1\}^4$.  Sign invariance makes this an integer
polynomial in its five arguments; every radicand has exponent at most
eight.  Our two equations are
\begin{align}
 P_R&=\delta^8\mathcal F_{16}
 \left(S(n+V)-A_c; r_a,r_b,r_e,\frac{S^2C_y}{\delta}\right)=0,
 \label{eq:s12-rational}\\
 P_I&=\mathcal F_{16}
 \left(S(n+\delta V)-A_c; r_a,r_b,r_e,S^2\delta C_y\right)=0.
 \label{eq:s12-integral}
\end{align}
The first formula means its symbolic polynomial after denominator
clearing, not numerical evaluation by division.  Both literal JSON
circuits contain only fixed integers, addition, subtraction, and
multiplication.

For completeness, the following polynomial recursion removes every
radical from the notation.  For a center $Z$, put $d=\delta$ and
$c_n=S^2C_y$ in~\eqref{eq:s12-rational}, or $d=1$ and
$c_n=S^2\delta C_y$ in~\eqref{eq:s12-integral}.  Define
\begin{align*}
 t_0&=Z^2+r_e,& t_1&=-2Z,& b_0&=r_a+r_b,\\
 q_0&=d(t_0-b_0)-c_n,&q_1&=dt_1,\\
 T_0&=d^2r_ar_b,&R_0&=dc_n,\\
 v_0&=q_0^2+r_eq_1^2-4(T_0+R_0b_0),&v_1&=2q_0q_1,\\
 z_{64}&=64T_0R_0,&&\\
 n_0&=v_0^2+r_ev_1^2-z_{64}t_0,&
 n_1&=2v_0v_1-z_{64}t_1.
\end{align*}
Then the polynomial output is $n_0^2-r_en_1^2$.  The independent oracle
checks all 495 coefficients against the sixteen-conjugate product.

\subsection{Unconditional integer zero-set bijections}

Write $P_R^{\rm old}$ and $P_I^{\rm old}$ for the pinned parent equations
with the same prefix and center $Sn-A_c$.  Every other abbreviation is
independent of $n$.  The maps
\[
 T_R:\ n_{\rm old}=n+V,\qquad
 T_I:\ n_{\rm old}=n+\delta V
\]
fix eleven coordinates.  Their inverse maps subtract $V$ and $\delta V$,
respectively.  These are mutually inverse polynomial maps of $\Z^{12}$,
without sign, nonzero, divisibility, or zero-equation hypotheses.
Direct substitution gives the identities
\[
 P_R=P_R^{\rm old}\circ T_R,\qquad
 P_I=P_I^{\rm old}\circ T_I.
\]
Consequently each pair has bijective complete integer zero sets, including
all exceptional branches as formal polynomials.  Composing these elementary
identities with the unchanged ordinary parent calibration gives, in the
same formalizability sense,
\[
 \mathrm{ZFC}\vdash
 \bigl(\forall x\in\Z^{12}\ P_R(x)\ne0\bigr)
 \leftrightarrow\operatorname{Con}(\mathrm{ZFC}),
\]
and the identical statement with $P_I$.  This composition introduces no
new native compiler claim.  The retained parent proof and its dependencies
remain part of the evidence; this section does not claim a new
machine-checked object-ZFC derivation.

\subsection{Exact degree certificates}

Only quantified coordinates contribute to total degree.  The prefix gives
\[
 \deg S=34,\quad\deg V=25,\quad\deg\delta=32,\quad
 \deg r_a=88,\quad\deg r_b=102,\quad\deg r_e=114,\quad
 \deg A_c=45.
\]
Put $z_0=Sn-A_c$, of degree 45.  For $P_R$, the relation
$C_y=(\delta+1)V^2-1$ gives the exact cancellation
\[
 \delta(SV+z_0)^2-S^2C_y
   =\delta z_0^2+2\delta SVz_0-S^2(V^2-1).
\]
The right side has degree at most 136.  Thus the norm register $q_0$ has
unique highest part $\delta_0(r_e)_0$ of degree 146, while
$(q_1)_0=-2\delta_0S_0V_0$ has degree 91.  Subsequent strict degree
comparisons give
\[
\begin{array}{c|rrrrrrrrr}
 &q_0&q_1&T_0&R_0&v_0&v_1&n_0&n_1&P_R\\\hline
 \deg&146&91&254&182&296&237&592&533&1184.
\end{array}
\]
Here and below a subscript zero on a polynomial denotes its highest
homogeneous component, independently of the temporary norm-register
names.  The full top form is
\[
 (P_R)_0=2^8 3^{24}S_0^{16}\delta_0^8V_0^8(D_g)_0^4.
\]
For $P_I$, the analogous identity is
\[
 (S\delta V+z_0)^2-S^2\delta C_y
 =z_0^2+2S\delta Vz_0-S^2\delta V^2+S^2\delta.
\]
Its unique highest term is $-S_0^2\delta_0V_0^2$, of degree 150.  The
remaining induction gives
\[
\begin{array}{c|rrrrrrrrr}
 &q_0&q_1&T_0&R_0&v_0&v_1&n_0&n_1&P_I\\\hline
 \deg&150&91&190&182&300&241&600&541&1200,
\end{array}
\]
and $(P_I)_0=S_0^{16}\delta_0^8V_0^{16}$.

These bounds are attained on the actual fixed tuple, not merely on generic
symbolic ports.  Selecting the largest $T$ power in each top form gives
\begin{align*}
 [j^{624}w^{144}v^{160}s^{64}r^{56}c^{112}h^{16}T^8]P_R
   &=2^{1256}3^{796}\beta^{624}>0,\\
 [j^{576}w^{144}v^{144}s^{80}r^{80}c^{160}T^{16}]P_I
   &=2^{1200}3^{720}\beta^{576}>0.
\end{align*}
The exponent sums are 1184 and 1200.  These are distinguished top-degree
monomials, not assertions that the full homogeneous components contain
only one term.  In particular the top part of $V$ contains
$4I_0T-R_0$, whose two summands both have degree five.  The exact checker
keeps the graded-lex leading term with $T$ first and confirms every
coefficient induction after the displayed cancellation.

\subsection{Cost and what the computations certify}

Under the same fully charged binary-operation convention as the parent,
the literal counts are
\begin{center}
\begin{tabular}{lrrrrr}
Equation&Multiplications&Additions&Subtractions&Total&Degree\\\hline
Old rational&93&26&26&145&1200\\
New $P_R$&93&27&26&146&1184\\
Old integral&89&26&26&141&1456\\
New $P_I$&90&27&26&143&1200
\end{tabular}
\end{center}
Thus $P_I$ reaches the old rational degree with two fewer operations,
while $P_R$ supplies a lower-degree option within this historical
comparison.  The saturated $123/1084$ equation of
Section~\ref{sec:integer12-saturated} improves both metrics of every row.

The standard-library verifier imports no new producer.  It proves the
whole-output substitution identities by a common expression DAG, checks
both integer inverse maps, binds all prefix formulas by independent
free-ring cuts, verifies every generic norm coefficient, and performs
the exact leading induction.  It also reconstructs all fixed coefficients
from the finite recipe and checks the live operation counts.  Normal and
optimized Python runs pass, and corrupted center and norm signs are
rejected.  These checks establish exact polynomial algebra and
transcription; the semantic conclusion uses the unbounded bijection proof
and the named ordinary parent theorem.

\section{An alternative source: an explicit equational proof kernel}
\label{sec:relational-kernel}

This appendix supplies a concrete local component of a different source
encoding.  It does not supply a new consistency equation.  The broad route
through finitely axiomatized class theory and relation algebra is prior
work: Carl and Moroz used finite NBG in their Diophantine provability
construction, and Formisano, Gavazzi, and Omodeo subsequently proposed a
Tarski--Givant relational version.  The latter paper does not print its
proper-axiom tuple; its Section~3.3 leaves those formulas outside its scope.
Consequently it cannot, by citation alone, calibrate a concrete new
polynomial to ZFC.  Our contribution here is an independently derived,
completely explicit local inference polynomial and an executable finite
proof checker for a supplied finite collection of equation schemata.
No priority claim is made for the local algebraic identity.

\subsection{A total coding of terms}

Work over $\mathbb N=\{0,1,2,\ldots\}$.  Let the abstract term language have
two constants $\epsilon,\iota$, two unary constructors $N,V$, and two binary
constructors $U,M$.  These will eventually stand for membership, identity,
complement, converse, union, and relation composition, respectively.  The
coding theorem itself is purely syntactic and does not require that
interpretation.  Set
\begin{align*}
 [\epsilon]&=0,& [\iota]&=1,\\
 [N(A)]&=4[A]+2,&[V(A)]&=4[A]+3,\\
 B(a,z)&=2(a+z)(a+z+1)+4z+4,\\
 [U(A,Z)]&=B([A],[Z]),&[M(A,Z)]&=B([A],[Z])+1.
\end{align*}
Every displayed constructor is a polynomial with integer coefficients.
In particular there is no division in the code of a compound term.
This is the polynomial form of the numbering in Section~3.1 of
Formisano--Gavazzi--Omodeo; the complete bijection proof is included here
to make the component independently checkable.

\paragraph{Theorem A.1 (total coding).}
The coding map is a bijection from these finite terms onto $\mathbb N$.

\paragraph{Proof.}
The Cantor function
\[
 C(a,z)=\frac{(a+z)(a+z+1)}2+z
\]
is a bijection $\mathbb N^2\longrightarrow\mathbb N$.  Indeed, for a given
$g$, the unique $s$ satisfying
$s(s+1)/2\leq g<(s+1)(s+2)/2$ gives
$z=g-s(s+1)/2$ and $a=s-z$.  Then $0\leq z\leq s$, so these are natural
numbers, and the formulas also prove uniqueness.  Moreover
$C(a,z)\geq\max(a,z)$.

The two constants occupy $0,1$.  Codes congruent to $2$ and $3$ modulo $4$
have the unique forms $4a+2$ and $4a+3$, respectively.  The remaining
numbers are $4(g+1)$ and $4(g+1)+1$ with $g\geq0$; the unique inverse of
$C$ identifies their two children.  Since $B=4(C+1)$, these are precisely
the two binary constructor cases.  In every case the child codes are
strictly smaller than the parent code.  Strong induction on a natural
number therefore constructs one finite term having that code; the
disjoint residue cases and the uniqueness of the Cantor inverse prove
that this term is unique.  This also proves injectivity of the original
coding map.  All inductions and integer computations are available in
ZFC. \hfill$\square$

Thus a natural-number witness never needs a separate certificate that
it denotes a valid term.  Equalities are stored as two coordinates
$(a,b)$, representing the equation between the uniquely decoded terms.
This avoids packing and immediately unpacking an equality as a further
Cantor pair.

\subsection{The seven-variable local polynomial}

Let $(a,b)$ and $(c,d)$ be two already proved equations and let $(u,v)$
be a proposed conclusion.  Introduce only one additional natural witness
$z$.  Define
\begin{align*}
 T&=(a-c)^2+(u-b)^2+(v-d)^2,\\
 q&=u-4a-2,\\
 A&=\bigl(q(q-1)\bigr)^2+\bigl(v-u-4(b-a)\bigr)^2,\\
 r&=u-B(a,z),\\
 D&=\bigl(r(r-1)\bigr)^2
       +\bigl(v-u-2(b-a)(a+b+2z+1)\bigr)^2,\\
 K(u,v,a,b,c,d,z)&=TAD.
\end{align*}
There are seven variables in $K$, not seven witnesses for the whole
provability problem: six coordinates are inputs and outputs of a single
local inference.  The history encoding is a separate cost.

\paragraph{Theorem A.2 (exact local inference).}
For natural values of its seven arguments, $K=0$ holds if and only if at
least one of the following five inference cases holds:
\begin{enumerate}
\item $a=c$, $u=b$, and $v=d$ (common-left transitivity);
\item $u=4a+2$ and $v=4b+2$ (complement congruence);
\item $u=4a+3$ and $v=4b+3$ (converse congruence);
\item $u=B(a,z)$ and $v=B(b,z)$ (union in the right argument);
\item $u=B(a,z)+1$ and $v=B(b,z)+1$ (composition in the right argument).
\end{enumerate}
Here ``in the right argument'' means adjoining the same right operand
$z$ to both members of a proved equality; it does not mean changing that
right operand.

\paragraph{Proof.}
The integer ring has no zero divisors, so $K=0$ is equivalent to the
vanishing of one factor.  Each factor is a sum of squares of integers.
For $T$, this gives exactly case~1.  For $A$, it gives
$q(q-1)=0$ and $v-u=4(b-a)$.  Hence $q$ is $0$ or $1$ and the same unary
tag is used on both sides, proving cases~2 and~3.

For $D$, it gives $r\in\{0,1\}$ and the stated difference equation.
Direct expansion gives the polynomial identity
\[
 B(b,z)-B(a,z)=2(b-a)(a+b+2z+1).
\]
Thus $u=B(a,z)+r$ and $v=B(b,z)+r$, proving cases~4 and~5.
Every listed case conversely annihilates its indicated factor.
The argument establishing this union of arithmetic cases works over all
integers; nonnegativity is used only for the term interpretation from
Theorem~A.1. \hfill$\square$

The two adjacent constructor tags are therefore selected without
introducing a Boolean variable.  The finite-difference identity also
avoids calculating $B(a,z)$ and $B(b,z)$ independently.  The included
straight-line circuit has 41 live binary additions, subtractions, and
multiplications; a square counts as one multiplication, and its constants
are $1,2,4$.  The circuit is an evaluator for this local polynomial, not
an operation count for a completed consistency equation.

\paragraph{Exact degree.}
The three factors have degrees $2,4,8$, respectively, so
$\deg K\leq14$.  Specializing $a=t$ and all other inputs to zero gives
leading terms $t^2$, $256t^4$, and $16t^8$.  Therefore the resulting
univariate polynomial has leading term $4096t^{14}$, proving
$\deg K=14$.

\subsection{Finite proof checking and the remaining obligations}

Fix a finite list of equation schemata.  Each side is a finite term
template over the six constructors, with a finite number of schematic
variables.  Supplying natural values for those variables produces an
exact instance by the constructor polynomials.  The code contains no
axiomhood oracle: the list of templates is literal input data.

A proof line records its conclusion, and either (i) a schema index and
its instantiation values, or (ii) two strictly earlier line indices and
the witness $z$.  A line of kind~(ii) is accepted exactly when $K=0$ for
its two parent equations.  One-parent congruence steps repeat the same
parent index.  The first line must be a schema instance.  This explicit
requirement excludes an unjustified cyclic first line.

\paragraph{Theorem A.3 (finite-checker adequacy).}
For the supplied schema list, the checker accepts a proof of a target
equation exactly when that equation has a finite derivation from those
schemata using the five inference cases in Theorem~A.2.

\paragraph{Proof.}
If the checker accepts, induction over its finite list proves each line.
At a schema line this follows by exact instantiation.  At an inference
line both parents have smaller indices, so the induction hypothesis and
Theorem~A.2 give its conclusion.  Conversely, list the nodes of a finite
derivation tree in a children-before-parent order.  Record the two
earlier indices for a transitivity step, duplicate the earlier index
for a congruence step, and use the code of its right operand for a binary
congruence.  For the other steps $z=0$ suffices.  Theorem~A.2 makes the
kernel vanish at each inference, and exact constructor evaluation
verifies every schema instance. \hfill$\square$

This theorem concerns the stated finite equational calculus, without a
claim that an arbitrary schema list expresses set theory.  A finished
replacement for the existing Con(ZFC) equation still requires all of:
\begin{enumerate}
\item A completely printed finite relational class-theory axiom tuple,
including the logical relation-algebra schemata and a precise target
equation expressing collapse.
\item Proof transformations, formalizable in ZFC, establishing that this
target is derivable exactly when the chosen ZFC calculus is inconsistent.
In particular the quasi-projection assumptions required for the
Tarski--Givant translation must be proved or included correctly.
\item An explicit polynomial compiler for arbitrarily long accepted
finite histories, with every bound, chronology condition, coefficient,
and witness accounted for.
\item Complete global variable, degree, coefficient, and evaluator counts
after that compiler is instantiated.
\end{enumerate}
None of these global costs can be inferred from the seven local
coordinates or 41 local operations.  The component provides a distinct
testable source direction, while leaving the established consistency
equation unchanged.

\subsection{An explicit combined row polynomial}

There is also a direct polynomial form of schema admission.  Suppose
the $j$th supplied schema has at most $m$ schematic variables.  Evaluating
its left and right templates using the constructor polynomials gives
fixed polynomials $L_j(\mathbf x),R_j(\mathbf x)\in
\mathbb Z[x_1,\ldots,x_m]$; unused coordinates are simply ignored.
Define
\[
 H(u,v,\mathbf x)=\prod_{j=1}^{s}
 \bigl((u-L_j(\mathbf x))^2+(v-R_j(\mathbf x))^2\bigr),
 \qquad
 F=K(u,v,a,b,c,d,z)H(u,v,\mathbf x).
\]
The empty product is $1$.  Over natural inputs, $F=0$ means exactly
that the output is a listed schema instance or follows from the two
parents by one of the five inference cases.  The same $m$ schema
coordinates may be shared between all $s$ factors: a zero product
requires only one schema factor to vanish.  Conversely witnesses for
that chosen schema can be padded with zeros.  This observation does
not share witnesses between multiple simultaneously required premises.

If the chosen logical basis includes reflexivity, histories can start
with the fixed justified row $(0,0)$.  Every later row can then carry two
strictly earlier parent indices, including axiom rows, which use the
initial row as harmless dummy parents.  Using $F=0$ for those later rows
avoids a separate Boolean axiom/inference tag.  The first row is fixed
and verified; it is never admitted merely by a self-referential
inference.  Induction gives the same soundness and completeness proof
as above.  This combined row polynomial still has an unbounded-history
compilation obligation.  In particular the variable count $7+m$ is a
local row count, not a variable bound for a Con(ZFC) equation.

\paragraph{Reproduction.}
Run \texttt{python3 relational\_kernel.py}.  The included finite checks
cover 10,000 complete code roundtrips, 10,000 Cantor inverse roundtrips,
2,048 binary constructor cases, 20,000 seeded arithmetic comparisons,
8,640 positive inference cases, and valid/invalid finite proofs.  These
checks supplement the universal proofs above; they do not establish the
missing NBG calibration or global history compilation.

\paragraph{Primary references.}
M.~Carl and B.~Z.~Moroz, \emph{On a Diophantine representation of the
predicate of provability}, MPI preprint~2006/129,
\url{https://webdoc.sub.gwdg.de/ebook/serien/e/mpi_mathematik/2006/129.pdf}.
A.~Formisano, I.~Gavazzi, and E.~G.~Omodeo,
\emph{Diophantine Modeling of Provability in Algebraic Logic}, CILC~2025,
\url{https://ceur-ws.org/Vol-4003/paper13.pdf}.


\begin{thebibliography}{9}\raggedright
\bibitem{repo} V. Reshetnikov, \emph{ProveIt: DiophantineConsistency},
snapshot \sha, inspected 24 October 2026.
\url{https://github.com/VladimirReshetnikov/ProveIt/tree/84e16e43092614bae262ca59651cb9d77e31a255/Computability/DiophantineConsistency}.
\bibitem{zero} ProveIt, \emph{Ten strictly positive integers, 99 operations,
degree 368}, \texttt{candidates/ten-positive-zero-offset}, at \cite{repo}.
The complete ordinary proof and literal parent are retained in this packet.
\bibitem{direct} ProveIt, \emph{Combined direct/zero-offset105/340},
\texttt{candidates/ten-positive-direct-zero-offset}, at \cite{repo}.
\bibitem{source61} ProveIt source61, \emph{Congruence carriers},
\texttt{con-zfc-diophantine-forms/61-congruence-carriers}; its rank,
index, and auxiliary completeness interfaces are restated in the retained
\texttt{ten-positive-scaled-norm.md} ordinary proof.
\bibitem{scalar12} ProveIt, \emph{Twelve unrestricted integers by a scalar carrier},
\texttt{candidates/scalar-carrier12}, at \cite{repo}.
The complete ordinary proof, signed-kernel interface, arbitrary-residue
interpolation theorem, and literal parents are retained in this packet.
\bibitem{dyadic} ProveIt incoming submission,
\emph{ConZFC Dyadic Pell Reductions, 23 October 2026},
\texttt{dyadic10/PROOF.md}.  Earlier distributed transport numerator and
separate upper-root centering; full archive at the pinned incoming snapshot.
\bibitem{carlmoroz} M. Carl and B. Z. Moroz,
\emph{On a Diophantine representation of the predicate of provability},
Max Planck Institute preprint 2006/129.
\url{https://webdoc.sub.gwdg.de/ebook/serien/e/mpi_mathematik/2006/129.pdf}.
\bibitem{fgo} A. Formisano, I. Gavazzi, and E. G. Omodeo,
\emph{Diophantine Modeling of Provability in Algebraic Logic},
CILC 2025, CEUR Workshop Proceedings 4003, paper13.
\url{https://ceur-ws.org/Vol-4003/paper13.pdf}.
\end{thebibliography}
\end{document}
